\documentclass[12pt,a4paper,openany,oneside]{book}

% --- Paquetes ---
\usepackage[utf8]{inputenc}
\usepackage[spanish, es-noshorthands]{babel}
\usepackage{amsmath, amsfonts, amssymb}
\usepackage{graphicx}
\usepackage{geometry}
\usepackage{hyperref}
\usepackage{booktabs}
\usepackage{fancyhdr}
\usepackage{listings}
\usepackage{xcolor}
\usepackage{tikz}
\usepackage{pgfplots}
\usetikzlibrary{arrows.meta, positioning, shapes.geometric}
\pgfplotsset{compat=1.18}

\geometry{margin=2.5cm}

% --- Colores y estilos para código Python ---
\definecolor{codegreen}{rgb}{0,0.6,0}
\definecolor{codegray}{rgb}{0.5,0.5,0.5}
\definecolor{codepurple}{rgb}{0.58,0,0.82}
\definecolor{backcolour}{rgb}{0.95,0.95,0.92}

\lstset{
    language=Python,
    backgroundcolor=\color{backcolour},   
    commentstyle=\color{codegreen},
    keywordstyle=\color{blue},
    numberstyle=\tiny\color{codegray},
    stringstyle=\color{codepurple},
    basicstyle=\ttfamily\small,
    breaklines=true,
    frame=single,
    showstringspaces=false
}

% --- Estilo de cabecera y pie ---
\pagestyle{fancy}
\fancyhf{}
\renewcommand{\headrulewidth}{0.4pt}
\renewcommand{\footrulewidth}{0.4pt}
\fancyhead[L]{\small Carlos César Sánchez Coronel}
\fancyhead[R]{\small Sesión 7: Regla de la Cadena y Derivadas de Orden Superior}
\fancyfoot[L]{\small Carlos César Sánchez Coronel}
\fancyfoot[C]{\thepage}

% --- Portada ---
\title{
    \textbf{Sesión 7} \\ 
    \Huge \textbf{REGLA DE LA CADENA Y DERIVADAS DE ORDEN SUPERIOR} \\
    \large La clave del backpropagation y la curvatura del error \\

}
\author{
    \small Por \\ \vspace{0.2cm}
    \Large \textsc{Carlos César Sánchez Coronel}
}
\date{2026}

\begin{document}

\maketitle
\tableofcontents
\addcontentsline{toc}{chapter}{Índice general}

% ------------------------------------------------------------------
% CAPÍTULO 7: SESIÓN 7 - REGLA DE LA CADENA Y DERIVADAS DE ORDEN SUPERIOR
% ------------------------------------------------------------------
\chapter{Regla de la Cadena y Derivadas de Orden Superior}

En esta sesión abordamos dos conceptos fundamentales del cálculo diferencial: la regla de la cadena, que nos permite derivar funciones compuestas, y las derivadas de orden superior, que nos informan sobre la curvatura de las funciones. Ambos son esenciales en inteligencia artificial: la regla de la cadena es el corazón del algoritmo de \textit{backpropagation} que entrena redes neuronales profundas, mientras que las derivadas de orden superior (en particular la matriz hessiana) se utilizan en optimizadores de segundo orden y en análisis de convergencia.

\section{Composición de funciones}

Antes de hablar de la regla de la cadena, recordemos qué es una función compuesta. Si tenemos dos funciones $f$ y $g$, la composición $f \circ g$ se define como:
\[
(f \circ g)(x) = f(g(x))
\]
Es decir, primero aplicamos $g$ a $x$, y luego aplicamos $f$ al resultado.

En inteligencia artificial, una red neuronal es una composición enorme de funciones. Por ejemplo, una red con una capa oculta puede expresarse como:
\[
\mathbf{y} = f_{\text{salida}}( \mathbf{W}_2 \cdot f_{\text{oculta}}( \mathbf{W}_1 \mathbf{x} + \mathbf{b}_1 ) + \mathbf{b}_2 )
\]
donde cada capa es una función lineal seguida de una no linealidad.

\subsection{Ejemplo sencillo}
Sean $g(x) = x^2 + 1$ y $f(u) = \sin(u)$. Entonces:
\[
h(x) = f(g(x)) = \sin(x^2 + 1)
\]
Para derivar $h$ necesitamos la regla de la cadena.

\section{La regla de la cadena}

La regla de la cadena establece que la derivada de una función compuesta es el producto de las derivadas de las funciones que la componen, evaluadas apropiadamente.

\subsection{Forma de Leibniz}
Si $y = f(u)$ y $u = g(x)$, entonces:
\[
\frac{dy}{dx} = \frac{dy}{du} \cdot \frac{du}{dx}
\]

\subsection{Forma de Lagrange}
Si $h(x) = f(g(x))$, entonces:
\[
h'(x) = f'(g(x)) \cdot g'(x)
\]

\subsection{Ejemplo paso a paso}
Derivar $h(x) = \sin(x^2 + 1)$.
\begin{align*}
\text{Paso 1: Identificar } u &= g(x) = x^2 + 1 \\
\text{Paso 2: } f(u) &= \sin(u) \\
\text{Paso 3: } f'(u) &= \cos(u) \\
\text{Paso 4: } g'(x) &= 2x \\
\text{Paso 5: } h'(x) &= f'(g(x)) \cdot g'(x) = \cos(x^2 + 1) \cdot 2x
\end{align*}

\subsection{Regla de la cadena generalizada}
Para composiciones múltiples, como $h(x) = f(g(h(x)))$, se aplica sucesivamente:
\[
h'(x) = f'(g(h(x))) \cdot g'(h(x)) \cdot h'(x)
\]
En general, la derivada es el producto de las derivadas de cada función en la cadena.

\subsection{Ejemplo con tres niveles}
Derivar $k(x) = e^{\cos(x^2)}$.
\begin{align*}
\text{Sea } a &= x^2, \quad b = \cos(a), \quad k = e^b \\
k'(x) &= \frac{dk}{db} \cdot \frac{db}{da} \cdot \frac{da}{dx} = e^b \cdot (-\sin(a)) \cdot (2x) \\
&= e^{\cos(x^2)} \cdot (-\sin(x^2)) \cdot 2x \\
&= -2x e^{\cos(x^2)} \sin(x^2)
\end{align*}

\section{Aplicación en IA: Backpropagation}

El algoritmo de \textit{backpropagation} (retropropagación) es una aplicación directa de la regla de la cadena para calcular gradientes en redes neuronales. La idea es: dado que la red es una función compuesta, la derivada de la función de pérdida respecto a los pesos de una capa se puede obtener multiplicando derivadas a lo largo de la red.

\subsection{Red neuronal simple: una neurona con dos capas}
Consideremos una red con una entrada $x$, una capa oculta con una neurona (activación sigmoide) y una capa de salida lineal (sin activación). La red calcula:
\begin{align*}
z &= w_1 x + b_1 \\
a &= \sigma(z) \\
\hat{y} &= w_2 a + b_2 \\
L &= \frac{1}{2}(\hat{y} - y)^2 \quad \text{(pérdida MSE)}
\end{align*}
Queremos $\frac{\partial L}{\partial w_1}$, $\frac{\partial L}{\partial w_2}$, etc.

\subsection{Derivadas parciales paso a paso}
Aplicamos la regla de la cadena:

\begin{enumerate}
    \item Primero, la derivada de $L$ respecto a $\hat{y}$:
    \[
    \frac{\partial L}{\partial \hat{y}} = \hat{y} - y
    \]
    \item Luego, $\frac{\partial L}{\partial w_2}$:
    \[
    \frac{\partial L}{\partial w_2} = \frac{\partial L}{\partial \hat{y}} \cdot \frac{\partial \hat{y}}{\partial w_2} = (\hat{y} - y) \cdot a
    \]
    Similarmente, $\frac{\partial L}{\partial b_2} = \hat{y} - y$.
    
    \item Para $w_1$, necesitamos más pasos:
    \[
    \frac{\partial L}{\partial w_1} = \frac{\partial L}{\partial \hat{y}} \cdot \frac{\partial \hat{y}}{\partial a} \cdot \frac{\partial a}{\partial z} \cdot \frac{\partial z}{\partial w_1}
    \]
    Calculamos cada factor:
    \begin{align*}
    \frac{\partial \hat{y}}{\partial a} &= w_2 \\
    \frac{\partial a}{\partial z} &= \sigma'(z) = a(1-a) \quad \text{(derivada de sigmoide)} \\
    \frac{\partial z}{\partial w_1} &= x
    \end{align*}
    Entonces:
    \[
    \frac{\partial L}{\partial w_1} = (\hat{y} - y) \cdot w_2 \cdot a(1-a) \cdot x
    \]
    y $\frac{\partial L}{\partial b_1} = (\hat{y} - y) \cdot w_2 \cdot a(1-a)$.
\end{enumerate}

Este proceso es exactamente lo que hace el backpropagation: propaga el error desde la salida hacia atrás, multiplicando por las derivadas locales.

\subsection{Visualización del flujo de gradientes}
\begin{center}
\begin{tikzpicture}[
    node distance=1.5cm,
    io/.style={circle, draw, minimum size=0.8cm},
    weight/.style={font=\small},
    >=Stealth
]
\node[io] (x) {$x$};
\node[io, right of=x] (z) {$z$};
\node[io, right of=z] (a) {$a$};
\node[io, right of=a] (yhat) {$\hat{y}$};
\node[io, right of=yhat] (L) {$L$};

\draw[->] (x) -- node[above] {$w_1$} (z);
\draw[->] (z) -- node[above] {$\sigma$} (a);
\draw[->] (a) -- node[above] {$w_2$} (yhat);
\draw[->] (yhat) -- node[above] {MSE} (L);

\draw[->, red, dashed] (L) -- node[below] {$\frac{\partial L}{\partial \hat{y}}$} (yhat);
\draw[->, red, dashed] (yhat) -- node[below] {$\frac{\partial \hat{y}}{\partial a}$} (a);
\draw[->, red, dashed] (a) -- node[below] {$\frac{\partial a}{\partial z}$} (z);
\draw[->, red, dashed] (z) -- node[below] {$\frac{\partial z}{\partial w_1}$} (x);
\end{tikzpicture}
\captionof{figure}{Flujo hacia adelante (negro) y hacia atrás (rojo) de gradientes.}
\end{center}

\section{Derivadas de orden superior}

La derivada segunda de una función mide la curvatura, es decir, cómo cambia la pendiente. En optimización, la segunda derivada nos dice si estamos en un mínimo, máximo o punto de silla.

\subsection{Definición}
Si $f$ es derivable, su segunda derivada es la derivada de la primera derivada:
\[
f''(x) = \frac{d}{dx} \left( \frac{df}{dx} \right) = \frac{d^2f}{dx^2}
\]

\subsection{Interpretación geométrica}
\begin{itemize}
    \item $f''(x) > 0$: la función es convexa (curva hacia arriba) en ese punto.
    \item $f''(x) < 0$: la función es cóncava (curva hacia abajo).
    \item $f''(x) = 0$: posible punto de inflexión.
\end{itemize}

\subsection{Ejemplo}
Sea $f(x) = x^3 - 3x$. Entonces:
\begin{align*}
f'(x) &= 3x^2 - 3 \\
f''(x) &= 6x
\end{align*}
En $x=0$, $f''(0)=0$ y hay un punto de inflexión. En $x=1$, $f''(1)=6>0$, mínimo local.

\subsection{Matriz Hessiana en varias variables}
Para funciones de varias variables, la segunda derivada se generaliza a la matriz Hessiana $\mathbf{H}$, cuyas entradas son:
\[
\mathbf{H}_{ij} = \frac{\partial^2 f}{\partial x_i \partial x_j}
\]
En optimización, el Hessiano nos da información sobre la curvatura de la función de coste. Los optimizadores de segundo orden (como Newton-Raphson) utilizan el Hessiano para acelerar la convergencia.

\subsection{Aplicación en IA: Análisis de la curvatura del error}
En problemas de machine learning, la función de pérdida es a menudo no convexa. Conocer la curvatura local (mediante el Hessiano) ayuda a determinar si un punto crítico es un mínimo, máximo o punto de silla. Los puntos de silla son comunes en altas dimensiones y pueden ralentizar el entrenamiento.

\section{Ejemplos con Python}

\subsection{Cálculo de derivadas de orden superior con SymPy}
\begin{lstlisting}[language=Python]
import sympy as sp

x = sp.Symbol('x')
f = x**3 * sp.exp(x) - sp.ln(x)

# Primera derivada
f1 = sp.diff(f, x)
print("f'(x) =", f1)

# Segunda derivada
f2 = sp.diff(f, x, 2)
print("f''(x) =", f2)

# Evaluar en un punto
valor = f2.subs(x, 1)
print("f''(1) =", valor.evalf())
\end{lstlisting}

\subsection{Derivadas parciales y Hessiano simbólico}
Para una función de dos variables:
\begin{lstlisting}[language=Python]
x, y = sp.symbols('x y')
f = x**2 * y + sp.sin(y)

# Derivadas parciales
df_dx = sp.diff(f, x)
df_dy = sp.diff(f, y)
print("df/dx =", df_dx)
print("df/dy =", df_dy)

# Hessiano (matriz de segundas derivadas)
hess = sp.hessian(f, (x, y))
print("Hessiano:\n", hess)
\end{lstlisting}

\subsection{Backpropagation manual en Python}
Implementemos el ejemplo de la red simple con una neurona oculta y una salida, usando derivadas calculadas simbólicamente y luego numéricamente.

\begin{lstlisting}[language=Python]
import numpy as np

# Datos de ejemplo
x = 2.0
y_true = 5.0

# Inicialización de pesos
w1, b1 = 1.0, 0.0
w2, b2 = 1.0, 0.0

# Forward pass
z = w1 * x + b1
a = 1 / (1 + np.exp(-z))  # sigmoide
y_pred = w2 * a + b2
loss = 0.5 * (y_pred - y_true)**2

print(f"Predicción: {y_pred:.4f}, Pérdida: {loss:.4f}")

# Backward pass (cálculo manual de gradientes)
dL_dy = y_pred - y_true
dy_dw2 = a
dy_db2 = 1.0
dy_da = w2

da_dz = a * (1 - a)
dz_dw1 = x
dz_db1 = 1.0

# Gradientes
grad_w2 = dL_dy * dy_dw2
grad_b2 = dL_dy * dy_db2
grad_w1 = dL_dy * dy_da * da_dz * dz_dw1
grad_b1 = dL_dy * dy_da * da_dz * dz_db1

print("Gradientes:")
print(f"  dL/dw1 = {grad_w1:.4f}")
print(f"  dL/db1 = {grad_b1:.4f}")
print(f"  dL/dw2 = {grad_w2:.4f}")
print(f"  dL/db2 = {grad_b2:.4f}")

# Actualización con descenso de gradiente (tasa de aprendizaje 0.1)
lr = 0.1
w1 -= lr * grad_w1
b1 -= lr * grad_b1
w2 -= lr * grad_w2
b2 -= lr * grad_b2

print("Pesos actualizados:")
print(f"  w1={w1:.4f}, b1={b1:.4f}, w2={w2:.4f}, b2={b2:.4f}")
\end{lstlisting}

\subsection{Verificación con diferenciación automática (usando JAX)}
Podemos comparar con JAX, que calcula gradientes automáticamente.

\begin{lstlisting}[language=Python]
import jax
import jax.numpy as jnp

def forward(params, x):
    w1, b1, w2, b2 = params
    z = w1 * x + b1
    a = 1 / (1 + jnp.exp(-z))
    y_pred = w2 * a + b2
    return y_pred

def loss(params, x, y_true):
    y_pred = forward(params, x)
    return 0.5 * (y_pred - y_true)**2

params = (1.0, 0.0, 1.0, 0.0)
x_val = 2.0
y_val = 5.0

grad_loss = jax.grad(loss)
grads = grad_loss(params, x_val, y_val)
print("Gradientes con JAX:", grads)
\end{lstlisting}

\section{Notación en papers científicos}
En artículos de deep learning, es común encontrar expresiones como:
\begin{itemize}
    \item Actualización de pesos mediante backpropagation:
    \[
    \frac{\partial \mathcal{L}}{\partial \mathbf{W}^{(l)}} = \delta^{(l)} (\mathbf{a}^{(l-1)})^\top
    \]
    donde $\delta^{(l)}$ es el error propagado hacia atrás.
    \item Hessiano de la pérdida: $\mathbf{H} = \nabla^2_{\mathbf{w}} \mathcal{L}(\mathbf{w})$.
    \item Aproximación de segundo orden en optimización:
    \[
    \mathbf{w}_{t+1} = \mathbf{w}_t - \eta \mathbf{H}^{-1} \nabla \mathcal{L}(\mathbf{w}_t)
    \]
    (método de Newton).
\end{itemize}

\section{Ejercicios propuestos}
\begin{enumerate}
    \item Deriva $f(x) = \ln(\sin(e^x))$ usando la regla de la cadena paso a paso.
    \item Calcula la segunda derivada de $g(x) = \frac{e^x}{x^2 + 1}$.
    \item Para la red neuronal del ejemplo, obtén las expresiones de las derivadas parciales si la función de activación fuera ReLU en lugar de sigmoide. ¿Qué cambios observas?
    \item Implementa en Python una red con dos neuronas en la capa oculta y calcula los gradientes manualmente (usa un loop sobre las neuronas).
    \item Investiga: ¿qué es el ``vanishing gradient'' y cómo se relaciona con la regla de la cadena?
\end{enumerate}

\section{Resumen}
En esta sesión hemos aprendido:
\begin{itemize}
    \item La regla de la cadena para derivar funciones compuestas, fundamental en backpropagation.
    \item Cómo se aplica la regla de la cadena para calcular gradientes en una red neuronal simple.
    \item Las derivadas de orden superior y su interpretación geométrica (curvatura).
    \item La matriz Hessiana y su papel en optimización de segundo orden.
    \item Implementaciones en Python con SymPy, NumPy y JAX.
\end{itemize}
Estos conceptos son la base matemática del entrenamiento de redes profundas y de métodos de optimización avanzados.

\end{document}